\documentclass[12pt]{article}
\usepackage[utf8]{inputenc}
\usepackage[T1]{fontenc}
\usepackage[margin=1in]{geometry}
\usepackage{setspace}
\usepackage{xcolor}
\usepackage[colorlinks=true, urlcolor=blue!60!black, linkcolor=blue!60!black]{hyperref}
\setlength{\parindent}{0em}
\setlength{\parskip}{0.6em}
\renewcommand{\baselinestretch}{1.0}
\begin{document}
\begin{center}
{\LARGE \textit{TransUnion LLC v. Ramirez}, 594 U.S. 413 (2021)\par}
\vspace{0.5em}
{\large Case Brief --- Two-Page Edition\par}
\end{center}

\section{Memory Jogger}
Only class members whose false ``terrorist'' credit alerts were actually sent to businesses had a concrete injury and standing to sue for damages (\href{https://supreme.justia.com/cases/federal/us/594/20-297/}{Justia}).

\section{Detailed Case Facts}
TransUnion's name-only screening falsely flagged consumers as potential matches to Treasury's terrorist list. A car dealer refused to sell Sergio Ramirez a car because of the alert, and TransUnion's follow-up mailings omitted required information. Of 8,185 class members, only 1,853 had reports sent to third parties.

\section{Procedural History}
A jury awarded the class over \$60 million under the Fair Credit Reporting Act; the Ninth Circuit found standing for all.

\section{Judicial Votes}
Reversed 5--4. Kavanaugh wrote for Roberts, Alito, Gorsuch, and Barrett; Thomas dissented with Breyer, Sotomayor, and Kagan; Kagan also dissented.

\section{Holding}
Every class member seeking damages must show a concrete harm bearing a ``close relationship'' to a traditionally recognized harm. Disseminated alerts resemble defamation; unsent internal files, unmaterialized risks, and harmless mailing defects do not.

\section{Analysis of Opinions}
\textbf{Majority:} Congress ``may not simply enact an injury into existence''; an undisclosed file is a defamatory letter in a desk drawer; risk of harm supports injunctions, not damages. \textbf{Thomas:} violation of a private statutory right has always sufficed; ``injury in fact'' dates only to 1970. \textbf{Kagan:} the Court made standing a tool of ``judicial aggrandizement.''

\section{Examples: Future Applications}
\textbf{Same-Side:} retaining former subscribers' data without disclosure confers no standing. \textbf{Opposite-Side:} disclosing viewing histories to an ad platform does. \textbf{Fence-Sitter:} sending debt data to a print vendor whose staff never reads it.

\section{Critique}
Originalists note that early Congresses authorized suits without proof of damage; progressives call the analogue test indeterminate and warn it pushes class actions into state courts.

\section{Key Quotations}
\begin{quote}
``Congress may enact legal prohibitions and obligations. And Congress may create causes of action for plaintiffs to sue defendants who violate those legal prohibitions or obligations. But under Article III, an injury in law is not an injury in fact. Only those plaintiffs who have been concretely harmed by a defendant's statutory violation may sue that private defendant over that violation in federal court.''
\end{quote}
\noindent --- Justice Kavanaugh, for the Court, 594 U.S. at 427.

\end{document}
